
\documentclass[12pt]{article}

%% Mathematics
%\usepackage{amsmath,amssymb,amsthm}
%\usepackage{amsmath}
\usepackage{amssymb} % \mathbb{}
\usepackage{mathrsfs}%

%% Physics
\usepackage{physics}

%% Customization of captions in floating environments
%\usepackage{caption} % for \caption*

%% Citation improvement
%\usepackage{cite}

%% Hypertext marks
\usepackage[hidelinks]{hyperref} % `hyperref` configuration must go before `orcidlink` to avoid the "Option clash for package hyperref" error.

%% ORCID
\usepackage{orcidlink} % to insert a hyperlinked ORCiD logo

%% Page output manipulation
%\usepackage[4]{pagesel}

%% Page layout settings
%\textwidth 160mm
%\textheight 230mm
%\voffset = -20mm

\title{Lattice gases: Grand partition function}

\author{R.V.~Romanik \orcidlink{0009-0008-2405-9038}
	\\ \small Yukhnovskii Institute for Condensed Matter Physics, NAS of Ukraine 
	\\ \small 1~Svientsitskii Street, 79011, Lviv, Ukraine 
	\\ \small romanik@icmp.lviv.ua}

%\date{} % to remove the date

\begin{document}
	
	\maketitle
	
	\begin{abstract}
		Abstract text.		
		\\[0.5cm]
		\textbf{Keywords:} Simple fluids, lattice gas.
	\end{abstract}
	
	\tableofcontents
	
	%\section{Introduction}
	
	%\pagebreak
	\section{General formulation for a many-particle system in a continuum}
	
	Consider a classical system of identical particles of mass $m$, confined to a macroscopic three-dimensional region $\mathcal{V} \subset \mathbb{R}^3$ of volume $V = \abs{\mathcal{V}}$, at temperature $T$ and chemical potential $\mu$. The grand partition function (GPF) is given by the following sum over all possible particle numbers $N$~\cite[Eq.~(2.4.6)]{HansenMcDonald13}: 
	\begin{equation}
		\label{def:gpf}
		\Xi(T, V, \mu) = \sum_{N=0}^{\infty}
		\frac{e^{N\beta\mu}}{h^{3N} N!} \int_{V^N} {\rm d} {\mathbf{r}}^N \int_{\mathbb{R}^{3N}} {\rm d} {\mathbf{p}}^N \exp\left[-\beta \mathcal{H}_N(\mathbf{r}^N, \mathbf{p}^N)\right] ,
	\end{equation}
	where $\beta = 1/(k_{\mathrm B} T)$ is the inverse temperature, $k_{\mathrm B}$ is the Boltzmann constant, $h$ is the Planck constant, and $\mathcal{H}_N$ is the Hamiltonian of the $N$-particle system. The positions and momenta of the particles are denoted by
	\begin{equation}
		\mathbf{r}^N \equiv \mathbf{r}_1, \ldots, \mathbf{r}_N, \quad \mathbf{p}^N \equiv \mathbf{p}_1, \ldots, \mathbf{p}_N,
	\end{equation}
	respectively, with the integration elements defined as
	\begin{equation}
		{\rm d} \mathbf{r}^N = \prod\limits_{i=1}^{N} {\rm d}^3 r_i, \quad {\rm d} {\mathbf{p}^N} = \prod\limits_{i=1}^{N} {\rm d}^3 p_i.
	\end{equation}
	
	The inclusion of the factor $1/h^{3N}$ in this definition ensures that $\Xi$ is dimensionless and consistent in form with the corresponding quantity of quantum statistical mechanics, while division by $N!$ ensures that microscopic states are correctly counted~\cite[p.~23]{HansenMcDonald13}.
	
	In the absence of external fields, the Hamiltonian is written as
	\begin{equation}
		\label{def:hamiltonian}
		\mathcal{H}(\mathbf{r}^N, \mathbf{p}^N) = K_N(\mathbf{p}^N) + W_N(\mathbf{r}^N),
	\end{equation}
	where
	\begin{equation}
		K_N(\mathbf{p}^N) = \sum_{i=1}^{N}\frac{\abs{{\mathbf{p}_i}}^2}{2m}
	\end{equation}
	is the kinetic energy and $W_N(\mathbf{r}^N)$ is the interparticle potential energy (the potential energy of interparticle interaction).
	
	Since the kinetic energy is quadratic in the momenta, the momentum integrations in Eq.~\eqref{def:gpf} can be evaluated analytically. Each momentum component contributes a factor $(2\pi m k_{\mathrm{B}} T)^{1/2}$. The GPF is therefore takes the form	
	\begin{equation}
		\Xi(T, V, \mu) = \sum_{N=0}^{\infty} \frac{\zeta^N}{N!} Z_N,
	\end{equation}
	where
	\begin{equation}
		\zeta = \frac{\exp(\beta \mu)}{\Lambda^3}
	\end{equation}
	is the activity, and
	\begin{equation}
		\Lambda = \frac{h}{\sqrt{2\pi m k_{\mathrm{B}} T}} = \left(\frac{2 \pi \beta \hbar^2}{m} \right)^{1/2},
	\end{equation}
	is the thermal de Broglie wavelength, with $\hbar = h/(2\pi)$ denoting the reduced Planck constant.
	
	The configuration integral $Z_N$ is defined as
	\begin{equation}
		\label{def:confInt}
		Z_N = \int_{V^N} \exp\left[-\beta W_{N}(\mathbf{r}^N)\right]{\rm d}\mathbf{r}^N,
	\end{equation}
	where the integration is performed over all configurations of $N$ particles in the volume $V$, with $Z_0 = 1$.
	
	It is common to assume that the interaction is pair-wise, and the potential energy of interaction can be written as
	\begin{equation}
		W_N(\mathbf{r}^N) = \frac{1}{2} \underset{i\neq j}{\sum_{i=1}^{N} \sum_{j = 1}^{N}} \Phi(\mathbf{r}_i, \mathbf{r}_j),
	\end{equation}
	where $\Phi(\mathbf{r}, \mathbf{r}')$ is the potential of interaction between a pair of particles located at $\mathbf{r}$ and $\mathbf{r}'$, respectively. The potential of interaction can, in turn, be divided into short-range and long-range parts. Conveniently, the short-range part is usually repulsive, while the long-range part is attractive~\cite{YRjps24}.

	%\pagebreak
	
	\section{General formulation for lattice gas models}
	In a lattice gas model consideration, the region $\mathcal{V}$ is partitioned into $\mathcal{N}$ non-overlapping congruent cubic cells $\Delta_l$, $l \in \{1, 2, \ldots, \mathcal{N}\}$, each of volume $v = \sigma^3$, where $\sigma$ is a linear length. The region $\mathcal{V}$ is therefore the union of all $\Delta_l$:
	\begin{equation}
		\mathcal{V} = \bigcup_{l=1}^{\mathcal{N}}\Delta_l,\qquad
		\Delta_l \cap \Delta_{l'} = \emptyset, \text{ if } l \neq l'.
	\end{equation}
	The volume of the whole system is thus $V = v\mathcal{N}$. 
	
	The interaction potential between a pair of particles now depends not on their particular positions $\mathbf{r}$ and $\mathbf{r}'$ or on their exact relative distance $\abs{\mathbf{r} - \mathbf{r}'}$, but rather on the particular cells that the particle coordinate belongs to:
	\begin{equation}
		\Phi(\mathbf{r}, \mathbf{r}') \to \Phi_{l, {l'}}, \text{ given } \mathbf{r} \in \Delta_{l}, \mathbf{r}' \in \Delta_{l'}.
	\end{equation}
	
	\textbf{Example of a nearest-neighbor attraction.} The most commonly considered interaction potential in the theory of lattice gases is the one allowing for at most single occupancy of cells and a constant attraction between particles occupying neighboring cells: 
	\begin{equation}
		\Phi(\mathbf{r}_i, \mathbf{r}_j) = 
		\begin{cases}
			+\infty, {\text{ if } \mathbf{r}_i \in \Delta_l, \mathbf{r}_j \in \Delta_l \text{ (i.e. two particles are in the same cell), }} \\
			-\varepsilon, \text{ if } \mathbf{r}_i \in \Delta_l, \mathbf{r}_j \in \Delta_{l'}, \text{ and } \Delta_l, \Delta_{l'} \text{ are the nearest neighbors},\\
			0, \text{ otherwise.}
		\end{cases}
	\end{equation}
	It can be expressed via the indicator function
	\begin{equation}
		\label{def:I}
		\mathbb{I}_{\Delta_l} (\vb{r}) = \left\{
		\begin{array}{ll}
			1, \quad \vb{r} \in \Delta_l,
			\\
			0, \quad \vb{r} \notin \Delta_l.
		\end{array}
		\right.
	\end{equation}
	Let us define the following auxiliary functions
	\begin{equation}
		S(\mathbf{r}, \mathbf{r}') = \sum_{l=1}^{\mathcal{N}} \mathbb{I}_{\Delta_l}(\mathbf{r}) \mathbb{I}_{\Delta_l}(\mathbf{r}'),
	\end{equation}
	and
	\begin{equation}
		A(\mathbf{r}, \mathbf{r}') = \sum_{\langle l, l' \rangle}
		\left[ \mathbb{I}_{\Delta_l}(\mathbf{r}) \mathbb{I}_{\Delta_{l'}}(\mathbf{r}') + \mathbb{I}_{\Delta_l}(\mathbf{r}') \mathbb{I}_{\Delta_{l'}}(\mathbf{r}) \right],
	\end{equation}
	where $\langle l, l' \rangle$ denotes an unordered pair of nearest-neighboring cells, counted once. Since every position belongs to exactly one cell, $S=1$ if the two particles are in the same cell and $S=0$ otherwise. Likewise, $A=1$ if their cells are nearest neighbors and $A=0$ otherwise.
	In terms of these functions, the interaction potential $\Phi$ becomes
	\begin{equation}
		\Phi(\mathbf{r}_i, \mathbf{r}_j) = 
		\begin{cases}
			+\infty, \quad S(\mathbf{r}_i, \mathbf{r}_j) = 1,\\
			-\varepsilon A(\mathbf{r}_i, \mathbf{r}_j), \quad S(\mathbf{r}_i, \mathbf{r}_j) = 0,
		\end{cases}
	\end{equation}
	or even can be written as a limit
	\begin{equation}
		\Phi(\mathbf{r}_i, \mathbf{r}_j) = \lim\limits_{K \to \infty} 
		\left[ K S(\mathbf{r}_i, \mathbf{r}_j) - \varepsilon A(\mathbf{r}_i, \mathbf{r}_j) \right].
	\end{equation}
	
	A yet another representation is in terms of occupation numbers $n_l$, which is the number of particles occupying a cell $\Delta_{l}$:
	\begin{equation}
		n_l = \sum_{i = 1}^N \mathbb{I}_{\Delta_{l}}(\mathbf{r}_i).
	\end{equation}
	If multiple occupancy is forbidden, then $n_l \in \{0, 1\}$, and 
	\begin{equation}
		\Phi(\mathbf{r}_i, \mathbf{r}_j) = -\varepsilon \sum_{\langle l, l'\rangle} n_l n_{l'}.
	\end{equation}
	
	In general, we can write the grand partition function for a lattice gas with unrestricted occupancy as
	\begin{eqnarray}
		\label{eq:Xi_lg}
		\Xi & = & \sum_{n_1 = 0}^\infty \left(\frac{v}{\Lambda^3}\right)^{n_1} \frac{{\rm e}^{\beta\mu n_1}}{n_1!} \ldots \sum_{n_{\mathcal{N}} = 0}^\infty \left(\frac{v}{\Lambda^3}\right)^{n_{\mathcal{N}}} \frac{{\rm e}^{\beta\mu n_{\mathcal{N}}}}{n_{\mathcal{N}}!} {\rm e}^{-\beta\mathcal{H}}.
	\end{eqnarray}
	
	A fundamental question arises whether a strict mathematical transformation is possible from the GPF represented in continuous coordinates, Eqs.~\eqref{def:gpf}--\eqref{def:confInt}, to the lattice gas representation~\eqref{eq:Xi_lg}. 
	
	One thing that seems to be obvious is that it is possible only if the interaction potential is discretized, i.e. the interaction depends on the cells where particles belong, but not on their precise coordinates. One striking example is an unrestricted occupancy model with the so-called Curie-Weiss interaction. The next section shows this transformation in details, based on the method of Ref.~\cite{KKD20}. Then we also consider the case of non-interacting particles for different maximum occupancies. We leave the question of a general case of generic discretised potential and of an arbitrary occupancy open.

	\section{Lattice gas with unrestricted occupancy and the Curie-Weiss interaction}
	
	%\pagebreak
	\section{Notes}
	In this work, we want to show how one can transform the GPF for a continuous system, Eq.~\eqref{def:gpf}, to that for a lattice gas system.
	
	When considering lattice gas systems, many authors neglect the kinetic energy part of the Hamiltonian~\eqref{def:hamiltonian}. We however want to study the lattice gas as a coarse-grained model of classical many particle systems in the continuum. One requirement to such a model is to grasp the liquid-gas phase transition.
	
	For a classical, single-occupancy lattice gas, the kinetic energy was considered in~Refs.~\cite{DM18,AM24}.
	
	The term `lattice gas' seems to be first coined by Lee and Yang in~\cite{YL52,LY52}. However, the model was considered by Lennard-Jones and Devonshire in Refs.~\cite{LD37,LD38}.
	
	Other references where a transition from a continuous system to a lattice system is considered: \cite{KMV03, LSS12, PBDGC26arxiv}
	
	Worth attention is Ref.~\cite{RB60}, where a cell model is considered in the grand canonical ensemble.
	
	
	%\section{Conclusions}
	
	% Bibliography
	\bibliographystyle{elsarticle-num}
	\bibliography{articles,books}
	
\end{document}
